\chapter{模型解释、误差分析与科学结论}
\chaptergoal{能够从总体指标进一步分析单个错误与子群体表现；正确使用模型系数、置换重要性等解释工具；辨别相关特征、分布变化与数据泄漏的影响；区分预测关系、相关性和物理机制；写出可复核且不过度外推的科学结论。}

机器学习项目结束时，最有价值的结果往往不是一个最大的准确率，而是清楚地知道模型在哪些条件下有效、为什么某些对象更困难，以及当前证据不能支持什么。本章把前面所有算法放在同一分析框架中：从总体评价回到错误样本，从局部表现回到数据与研究问题。
\section{从单个指标到错误样本}

\textbf{错误分析}是把汇总的评价拆回具体记录或有意义的分组，寻找重复出现的失败情形。例如人工设定一个100条记录的任务，常见类有90条且全部正确，稀有类有10条却只识别对2条；总体准确率为$(90+2)/100=92\%$，稀有类召回率却只有20\%。若研究重点恰好是稀有类，92\%并不能回答任务是否成功。

这里的分组应在开发分析中根据任务确定，并同时报告每组样本数。只有两条记录的组出现50\%正确率，与有两千条记录的组出现50\%正确率，所支持的稳定性判断不同。

一个模型的平均准确率或 MAE 只是对许多不同记录的汇总。如果少数科学上重要的对象被大量误判，总体指标仍可能很好看。分类问题应该检查混淆矩阵和逐类别召回率；回归问题应该检查残差与目标值、输入特征、观测条件和对象群体的关系。将错误记录按原因整理，有助于形成真正可改进的研究问题。

建立错误样本表时，应保存唯一标识、真实目标、预测结果、相关输入、质量字段和模型版本。对天体分类，某些误判可能集中在暗弱目标或颜色重叠区域；对物理实验，某些残差可能集中在高频或特殊流速条件。诊断要同时查看样本数量，避免只凭极少数记录的平均值作出稳健性结论。

\fig{ch16_wine_confusion.pdf}{Wine 逻辑回归开发数据上的示例混淆矩阵，演示逐类别错误分析。具体结果来自图形脚本的固定划分。}

\subsection{把一次错误追溯到它所在的样本范围}
整体分数把每条记录压缩到一个总数，适合比较方案，却可能掩盖局部失败。例如一万个常见对象和十个罕见对象混合时，模型在罕见对象上全部失败，只会轻微改变总体准确率。科学应用往往更关心这些对象，因此报告需要同时列出总体、各类或各条件的记录数与误差。

错误表应保存对象标识，而不是只保存重排后的行号。一次预测应能对应回原始输入、目标来源、质量字段、训练划分和模型版本。若只保留“第17行预测错了”，而Notebook前面又筛选或排序，后续很难知道真正讨论的是哪个对象。

\worked{总体均值相同，失败方式可以完全不同}
两组各4条回归残差，A为 $(-1,-1,1,1)$，B为 $(0,0,0,2)$。两组RMSE都为1，但A没有平均偏差、每条都差1；B只有一条大误差，平均偏差为0.5。若任务对偶尔的大误差极敏感，二者风险不同；若任务关注系统偏移，也不同。不能用相同RMSE推断模型行为等价。

这个例子提示至少查看残差均值、绝对误差分位数、尾部记录和输入范围。每项统计回答一个问题，不必无目的地罗列几十个指标，但应让指标服务于预先定义的科学或工程需求。

\subsection{先检查数据与评价，再解释模型}
发现大误差时，首先核对字段单位、目标与输入的对应、缺失处理和对象身份。一次错误连接表就可能让所有模型都变差；把摄氏度与开尔文混合，也可能形成看似复杂的分组。确认这些基础信息后，再检查模型表达与训练配置。科学解释应从可排查的计算事实逐步推进，而不是首先假设发现新物理。

\section{分类错误与回归残差}
分类模型可以把预测概率、真实类别和最终类别并排分析。需要区分预测概率接近阈值造成的边界错误与概率非常极端的错误。对于稀有类，错过一个真正候选与误报一个背景对象的代价可能不同。若实验需要后续人工检查，候选排序和可处理数量也应进入评价设计。

回归常用残差 $e_i=\hat y_i-y_i$，本书保持预测减去目标的符号约定。残差随输入有系统趋势，可能提示遗漏特征或模型形式不足；不同工况的残差方差显著不同，可能提示测量误差变化、样本不足或物理机制随条件变化。残差分析的第一步是描述现象，第二步才是提出需要额外证据检验的原因。

\subsection{分组残差怎样从数据中计算}
对连续输入，可以预先按有意义区间分组，再在每组报告样本数、平均残差和RMSE。以共同的人工响应验证数据为例，随机森林的60条验证记录按主要输入 $x$ 分成三组：
\begin{center}\small
\begin{tabular}{lrrr}\toprule
输入范围&记录数&平均预测减目标&RMSE\\\midrule
$[-2.5,-1)$&18&0.0912&0.3383\\
$[-1,1)$&27&$-0.0551$&0.3277\\
$[1,2.5]$&15&$-0.0470$&0.3806\\\bottomrule
\end{tabular}\end{center}
\fig{v3_ch16_residuals.pdf}{人工响应验证记录的预测残差。分组汇总与逐点图共同用于检查偏差和局部散布，不从一组较大RMSE直接推断物理原因。}

右侧区间误差较大，但它只有15条记录，不能把几百分之一的差异直接宣布为稳定规律。下一步可以检查该区域训练覆盖、目标变化和个别大残差；若据此修改模型，这些验证记录继续属于开发过程。最终测试一旦被反复用于设计新的分组和筛选条件，就不能继续被称为完全独立确认。

\subsection{测量噪声与模型残差不是同一个量}
假设观测 $y=f(\mathbf x)+\epsilon$。残差包含模型对 $f$ 的误差，也包含观测噪声 $\epsilon$；若输入本身有误差，还会增加另一类影响。即使模型完全恢复真实平均关系，单次观测残差也不必为0。因此训练误差逼近0可能是在拟合噪声，而非发现更精确的物理关系。

若每条观测提供可信的目标标准误差 $\sigma_i$，可考察归一化残差 $(y_i-\hat y_i)/\sigma_i$，但它只有在误差模型、相关性和预测不确定性处理合适时才具有预期统计解释。若预测本身也不确定，分母只使用测量误差会夸大偏离；若误差之间相关，仅逐条除以标准差也不能消除相关结构。

\subsection{概率错误与类别错误要分开}
一个分类器把100条记录都预测为正确类别，但每次真实类别概率都只有0.51；另一个在99条上概率0.99且正确，在1条上非常自信地错误。准确率、交叉熵和应用风险可能给出不同偏好。概率模型应检查不仅“分对多少”，还包括概率是否与长期频率相符。

例如把预测概率在0.7至0.8的对象分为一组，若有100个、其中正类60个，则这组平均概率约0.75高于实际比例0.60，提示过度自信。但有限样本、分组规则和总体变化都会影响判断。可靠性图是诊断工具，不是单独的因果或校准保证；重新校准需使用开发或独立校准数据。

\section{特征系数与置换重要性}

\textbf{模型解释}研究模型怎样利用输入形成输出。解释单条预测称为局部解释，概括一组记录的模型行为称为总体解释；二者都不自动解释真实世界的因果机制。

\textbf{置换}在这里指打乱一列数值的行对应关系，保持这一列已有取值，却破坏它与其他列、目标之间原来的匹配。例如一列 $[1,2,3,4,5]$ 被改排为 $[3,5,1,2,4]$，其他列不动。保持模型不重新训练，再比较验证表现。若原准确率为0.90，置换后为0.70，则按准确率下降定义的重要性为0.20；重复几次置换可观察偶然重排带来的波动。对MAE这样的误差指标，应比较误差增加量并说明方向，不能不看符号就排序。

线性模型系数可以描述在其他输入保持不变时，改变一个特征对模型预测的影响；但不同单位、标准化和强相关特征会影响解释。树的内置重要性反映模型内部划分或损失减少的某些统计量，并不天然等于科学变量的重要性。

置换重要性（permutation importance）通过打乱某个特征列，再观察预测评价指标的变化，估计当前模型在当前数据范围对该列信息的依赖程度。它应在用于解释的独立或验证数据上进行，并确保没有把最终测试集变成反复筛选特征的工具。对于高度相关的特征，打乱一个变量之后，其他变量可能保留相近信息，因此单列重要性可能偏低；反过来，置换产生不自然的特征组合也可能改变解释。

\begin{lstlisting}[caption={Wine 分类模型的置换重要性示例}]
from sklearn.datasets import load_wine
from sklearn.model_selection import train_test_split
from sklearn.pipeline import make_pipeline
from sklearn.preprocessing import StandardScaler
from sklearn.linear_model import LogisticRegression
from sklearn.inspection import permutation_importance

B = load_wine(as_frame=True)
X_train, X_val, y_train, y_val = train_test_split(
    B.data, B.target, stratify=B.target,
    test_size=0.3, random_state=42)
model = make_pipeline(StandardScaler(), LogisticRegression(max_iter=2000))
model.fit(X_train, y_train)
r = permutation_importance(model, X_val, y_val,
                           n_repeats=10, random_state=42,
                           scoring="accuracy")
for name, value in sorted(zip(B.feature_names, r.importances_mean),
                          key=lambda z: -z[1])[:5]:
    print(name, value)
\end{lstlisting}

\fig{ch16_permutation.pdf}{Wine 分类任务中，在验证数据上计算的置换重要性及重复置换波动。这是模型依赖性诊断，不是变量的因果贡献。}

\subsection{把系数放回它的条件含义}
在线性回归中，一个系数表示其他模型输入保持不变时，当前输入变化一单位所对应的预测变化。若输入高度相关，这种“其他不变”的组合可能在实际数据中很少出现；系数依然是数学模型的一部分，但对现实干预的解释需要额外假设。

标准化系数表示输入变化一个训练标准差的影响，原单位系数表示一个物理单位的影响。直接比较厘米特征与米特征的原始系数大小没有意义。对于非线性特征展开，原变量变化可能同时改变一次项、平方项和交互项，不能只读取其中一个系数。

\subsection{置换实验要保持哪些东西不变}
固定已训练模型和一批开发验证记录，计算基准分数 $S$。只打乱某一特征列在样本之间的对应，其他列与目标保持原位，再计算 $S_j^\mathrm{perm}$。若分数越高越好，重要性为 $S-S_j^\mathrm{perm}$；若使用误差，则可定义为置换后误差减原误差。务必先确认方向，不把负RMSE的变化再次随意取负。

重复置换获得波动，是对列打乱这一随机操作的条件变异，不包含重新采集训练数据、重训模型或未来目标噪声。某列置换后表现偶然改善，可能来自有限样本、模型误用该列或相关性，不应强行截成0后宣称它完全无关。

\worked{重复特征使“逐列贡献”变得不唯一}
假设两列都准确记录同一个变量，模型可以选择任一列，或把系数分摊。若树只使用第一列，第二列置换的重要性可能为0；若在重新训练时移除第一列，第二列可以替代它，性能仍很好。于是“固定模型对某列的依赖”“删列后重新训练的信息损失”和“变量在科学上是否重要”是三个不同问题。

可以将相关列一起置换，检查一组信息的共同贡献；也可做有控制的删列重训比较。无论采取哪种，都应说明具体问题和训练成本，不能把多个不同定义的重要性当作一个唯一真实排序。

\section{相关特征和解释的不确定性}

\textbf{选择效应}指进入数据表的对象不是目标总体的随机代表，而受观测、实验或整理条件筛选。例如只给明亮天体取得可靠类别标签，训练表中明亮对象就比实际应用总体更常见。\textbf{代理信息}是能帮助预测却不是所研究物理机制本身的线索，例如某类对象恰好只由一台设备测量，模型可能利用设备差异代替类别本身的规律。理解这些术语之后，才能讨论特征重要性为何会随数据范围变化。

若两列输入几乎表达同一信息，模型可能在二者之间分配权重，或者只使用其中一个。删去某列后整体表现仍然良好，并不证明该变量与物理任务无关；它可能被其他变量替代。建议通过相关矩阵、成组置换、删除特征后重新训练等对照实验，检查结论是否依赖特定特征组合。

解释还受样本选择影响。如果某些类别只在高信噪比区域拥有标签，模型在这个子样本中发现的重要特征关系，未必适用于全体观测目标。模型可能利用与研究目标相关但属于测量条件的代理信息。这些依赖有时在当前数据上提升预测，却会在仪器或处理方式改变时失效。

\subsection{重采样需要说明哪些量固定、哪些量重新估计}
对一组已经固定的验证预测做行Bootstrap，每次有放回抽取验证记录并重新计算RMSE，得到的是在当前模型、当前经验验证分布条件下的指标波动。如果每次还重新抽取训练记录、拟合预处理和模型，则会额外包含训练过程变化。若还在每次重采样中重新选择超参数，又在评价整个选择流程。三者回答的问题不同，计算代价也不同。

本轮对60条人工响应验证记录及固定随机森林预测进行1000次成对行重采样，RMSE为0.3448，Bootstrap分布的2.5\%与97.5\%分位数约为0.2761和0.4248。这个范围描述该条件下的指标重采样分布，不是每个新对象的预测区间，也不是已经涵盖仪器变化的总体保证。

若一组记录来自同一对象或实验批次，按行重采样会把相关观测当成独立。应根据目标以对象或批次为单位抽样，并整体保留组内记录。重复观测数量多的对象也可能在按行指标中占更高权重，报告需说明是按记录还是按对象平均。

\subsection{平均响应区间与新观测区间}
在 $y=f(x)+\epsilon$ 的框架中，估计平均响应 $f(x)$ 的不确定性与预测下一次 $y$ 的不确定性不同。后者还包含新观测的噪声。在简单独立情形，预测误差方差可以包含参数估计贡献和噪声贡献；只观察多次重训模型的预测差异，通常没有包含全部未来噪声。

训练样本增加可以减少某些估计不确定性，却未必减少传感器本身的测量噪声。若新的观测条件完全不在训练范围内，熟悉分布下的区间构造也可能失效。附录E介绍高斯过程时，会在明确概率假设下区分这些量；本章先建立不同问题不能互换的认识。

\subsection{特征解释的不稳定性也应作为结果}
若多次开发重采样的系数符号变化，而预测误差差不多，可能存在共线性或多种近似等效表达。此时报告一个固定符号的物理结论并不稳健。可以展示系数分布、相关矩阵、组合特征的共同效果，或者说明当前数据无法唯一辨识各变量作用，而不是只选择符合预期的一次模型。

\section{对象群体、实验条件与分布变化}

\textbf{分布变化}指模型实际接收的数据与训练数据在取值范围、比例或输入--目标关系上有所不同。它不只是多了一些新行。例如原训练数据主要是高信噪比测量，应用数据多数信号较弱，即使字段名完全一致，输入误差和目标关系也可能变化。\textbf{信噪比}描述信号相对噪声的强弱，具体计算方式需以数据源定义为准。

仅重新计算应用数据的均值和标准差并不能保证解决分布变化，甚至可能破坏已训练模型的输入含义。先检查差异，再决定是否需要新标签、重新训练或限制应用范围。

机器学习模型的泛化需要指定目标群体。对 SDSS 测光分类，按星等、颜色、信噪比或天区分组检查错误；对 Airfoil，按频率、流速和攻角检验残差；对材料预测，按化学组成或材料族讨论新组合。每组样本数不同，指标波动的解释也不同。对于新仪器、新观测夜或新实验工况，模型可能遭遇训练时没有覆盖的分布。

\fig{ch16_group_residual.pdf}{使用 sklearn 随附的 Diabetes 回归示例作分组残差诊断的演示；真正的天文与物理结论须在相应专业数据上重新计算。}

对分布变化的检查可以从训练和应用数据的特征分布比较开始，但即使输入边缘分布相近，特征与目标关系也可能改变。应记录模型支持的样本范围，并在部署前定义何时需要重新训练、校准或拒绝给出确定预测。

\subsection{几种分布变化需要不同的检查}
输入分布变化指应用记录的特征范围或比例与训练不同，例如新天区更暗弱、实验温度更高。类别先验变化指不同类别的总体比例改变，即使条件特征分布近似相同，候选精确率也会变化。输入与目标的关系变化则更直接，例如仪器老化使相同输入对应不同输出。它们可能同时发生，但不能都用同一个随机交叉验证数字代表。

可以分别构造按对象、时间、地域、材料族或实验工况留出的开发检验。选择依据是未来要面对的变化，而不是哪一种划分给出最高分。若真正部署到新仪器，但数据中没有多仪器记录，就只能承认该问题尚未检验；不能把随机划分换名为跨仪器泛化。

\subsection{覆盖范围检查如何进入结果表}
对每条待预测记录，保存其关键输入是否超出训练范围、与近邻训练样本的距离、质量标志和缺失模式。这些不是普适的可靠性概率，但能帮助识别明显缺乏支持的输入。单列范围都正常也不保证联合组合正常，例如训练只见过高温高压和低温低压，后来遇到高温低压，可能仍是新的组合。

外推检验最好专门设计，而不是仅从随机保留数据中寻找几个极值。沿时间的预测应尊重先后顺序，组合信息可能需要避免把未来统计量提前用于预处理。对材料族外推，要检查化学组成是否以不同命名重复出现在两侧。

\subsection{评价应允许“暂不判断”}
某些科学任务可以设置人工复核或拒绝区间：模型对覆盖不足、质量低或概率难以区分的对象不直接给最终身份。此时需要同时报告覆盖率和被覆盖对象的错误率，不能只留下容易样本后报告高准确率而不说明删去了多少对象。拒绝本身也是一个决策，需要在开发阶段规定并检验。

\section{预测关系与科学解释}

\textbf{预测关系}指利用已有信息估计未知目标的有效联系；\textbf{统计关联}指变量在数据中共同变化的模式；\textbf{因果作用}则要求讨论改变某个因素会怎样影响结果。三者需要的证据不同。例如两列都随同一个温度条件变化，即使互相很容易预测，也不代表改变其中一列就能使另一列跟着改变。

\textbf{不确定性}也要说明对象。测量误差来自观测或实验过程；重复训练的波动来自数据抽样或初始化；交叉验证的折间分数波动来自几次不同训练/评价划分；针对一个新对象的预测区间则要描述该对象目标可能落入的范围。它们不能用同一个“误差条”不加说明地替代。

相关与预测不能直接推导出因果机制。例如，某个测光颜色是天体类别的有效预测变量，并不证明该颜色单独决定了物理类别；树分裂阈值也不等于物理临界常数。科学解释需要结合已有理论、独立观测、受控实验或明确的因果识别设计。

模型预测的“重要性”可能只说明当前算法依赖某个输入；模型在未见数据上的“泛化”只说明某种评价范围内的预测能力；而物理机制需要不同性质的证据。写作中应分别使用“预测”“相关”“一致于某种解释”和“支持特定机制”等表述，不要把它们混为同一强度的结论。

还应区分交叉验证折间波动、重新训练的随机性、测量误差和单个预测值的不确定性。它们来源不同，不能把折间标准差直接写成每颗天体的预测置信区间。更完整的不确定性模型和高斯过程可在附录专题中继续学习。

\subsection{从关联到机制还缺少什么}
若模型利用某个测光颜色预测天体类别，说明该颜色在当前样本和其他输入条件下提供了预测信息。它可能与温度、红移、消光或仪器选择共同相关。要把它解释成单一机制，需要结合理论、独立测量或受控变化；预测表现本身不能排除这些替代解释。

物理实验中的受控设计可以提供更强证据。例如只改变一个输入并保持其余条件，观察响应；但现实数据常不是这样生成的。机器学习可以帮助建模、发现残差结构和提出假设，却不能自动把观察性目录转变为随机干预实验。

\worked{一个看起来重要、却并非原因的编号}
假设实验先在低温条件测量批次1至50，再在高温条件测量51至100，目标也随温度变化。批次编号可以很准确地预测目标，但它并不制造物理响应。如果未来重新排列测量顺序，编号规则可能失效。模型重要性指出它利用了什么，不直接回答对象为什么这样变化。

这不意味着时间或编号永远不能作为输入。若任务就是预测同一过程随时间的变化，时间可以有意义。关键是输入是否具有明确可部署含义，以及验证设计是否检验了预期应用，而不是简单把某类字段永久列为禁用。

\subsection{将结论分成三个层次}
第一层是计算结论：在指定数据和划分下，某模型达到多少误差。第二层是经验关系：哪些对象或条件下预测较好，哪些关联在多种检验中保持。第三层是物理机制或因果解释：需要哪些额外理论和证据。报告可以到达不同层次，但每一步都要有对应支持。

专业教材应训练学生把“目前只能支持第一层”说清楚，而不是为了让报告显得深刻就跃迁到第三层。一个明确指出模型失败条件的结果，可能比没有限制的高分宣称更有科研价值。

\section{结果报告与可复现记录}
一份合格的科学机器学习报告至少说明：问题及目标应用、样本与来源、字段与质量要求、数据划分、基线与候选模型、评价指标、关键结果、错误和分组分析、局限、代码与软件版本。图表应标明样本范围、变量名称与单位，以及是否来自训练、验证或最终测试。

推荐在项目中为每一张表和图保留生成脚本，并记录源数据的唯一版本。结果涉及随机过程时至少报告随机种子和主要配置；若表现对初始状态明显敏感，应给出波动范围。模型选择后保留一次独立测试的原始记录，避免因为后期改图、改表或调参而造成结果口径漂移。

\subsection{让每个结论找到对应证据}
建议结果段落围绕一个具体问题组织：先说明本次比较的对象和评价方案，再报告主要数值及波动，接着指出图中的关键行为，最后解释当前证据支持的结论。例如“在按批次留出的开发折上，二次模型相对线性基线减小了高输入区的系统残差”，比单独写“效果显著提升”更可核查。

表格标题应写清数据范围、单位、训练或验证身份，以及平均方式。图中坐标、图例和样本数不能依赖教师口头补充。同一章若同时出现人工响应、住房和真实天体，应在每张图中标明对象，避免读者把机制演示的数值当成观测结果。

\subsection{可复现记录的最小结构}
保存原始来源和校验值、筛选代码、字段字典、对象标识、划分索引、软件版本、模型配置、预测表和绘图脚本。只保存最终模型文件不足以复现训练，只保存Notebook截图也不能重新计算。涉及远程数据时应记录取得日期和版本；若本次没有取得该数据，就把结果状态明确写成未执行。

同一环境重现到小数点后一致是数值复现；在其他种子、对象划分和合理配置下保持结论，是稳定性。二者都重要，但一个哈希或固定种子不能替代另一个。课程报告应同时回答“别人怎样重做”和“结论对合理变化是否敏感”。

\chapterpractice{从 Wine 或住房任务选一个已训练的模型，分析错误样本或残差；使用置换重要性检查特征依赖，并说明相关特征造成的解释限制；在一个专业案例上按亮度、实验条件或对象群体进行分组评价。}
\chapterexercise
\begin{enumerate}
\item 为什么一个总体准确率很高的模型仍可能对稀有科学类别没有实用价值？
\item 某个变量的置换重要性为零，是否说明该变量没有物理作用？
\item 折间评价标准差、测量误差和预测区间分别在描述什么？
\item 说明训练数据范围外的模型输出应如何在报告中表述。
\item 设计一张模型比较表与一幅分组误差图，并逐项写明应当包含的元数据。
\end{enumerate}

\section*{深化练习与方法检验}
本组题强调解释、计算与验证；D编号在本章内独立计数。计算题应写出所用定义和中间步骤，实验题保留实际输出，包括不支持预期的结果。
\begin{enumerate}[label=D\arabic*.]
\item 对已经固定的模型和一张验证预测表做bootstrap，得到RMSE的百分位区间。这个区间是否包括重新采集训练集、重新调参和改变仪器的全部不确定性？
\item 模型使用两列高度相关的测量，逐列置换重要性均很小。能否说两个量都无关？设计联合置换或重新训练实验，并说明其解释对象。
\item 一个模型在低信噪比样本的RMSE大于高信噪比样本。至少提出三种可能原因，并说明仅根据这一张图不能区分什么。
\item 假设校准残差绝对值排序为 $r_{(1)}\le\cdots\le r_{(19)}$。按正文简单分裂保形方法构造90\%预测区间时应取哪个秩？指出覆盖结论需要的关键条件。
\end{enumerate}
